% Mechanistic Validity — NEMI 2026 poster, v1
% tikzposter 3x3 block format, matching grokking_poster_v2.tex style.
%
%   lualatex mechval_poster_grok_v1.tex

\documentclass[25pt, a0paper, portrait]{tikzposter}

\usepackage{amsmath,amssymb}
\usepackage{booktabs}
\usepackage{graphicx}
\usepackage{enumitem}
\usepackage{microtype}

% ── Color palette ───────────────────────────────────────
\definecolor{confirmed}{HTML}{2980B9}
\definecolor{partial}{HTML}{E67E22}
\definecolor{inconclusive}{HTML}{95A5A6}
\definecolor{untested}{HTML}{BDC3C7}
\definecolor{disconfirmed}{HTML}{E74C3C}
\definecolor{darktext}{HTML}{2C3E50}
\definecolor{blockbg}{HTML}{FFFFFF}

% ── Theme ───────────────────────────────────────────────
\usetheme{Default}

\definecolorstyle{MechValColors}{
  \definecolor{colorOne}{HTML}{2C3E50}
  \definecolor{colorTwo}{HTML}{3498DB}
  \definecolor{colorThree}{HTML}{ECF0F1}
}{
  \colorlet{backgroundcolor}{colorThree}
  \colorlet{framecolor}{colorOne}
  \colorlet{titlefgcolor}{white}
  \colorlet{titlebgcolor}{colorOne}
  \colorlet{blocktitlebgcolor}{colorTwo}
  \colorlet{blocktitlefgcolor}{white}
  \colorlet{blockbodybgcolor}{blockbg}
  \colorlet{blockbodyfgcolor}{darktext}
  \colorlet{notefgcolor}{darktext}
  \colorlet{notebgcolor}{partial!15}
  \colorlet{notefrcolor}{partial}
}
\usecolorstyle{MechValColors}

\defineblockstyle{MechValBlock}{
  titlewidthscale=1.0, bodywidthscale=1.0, titleleft,
  titleoffsetx=0pt, titleoffsety=0pt,
  bodyoffsetx=0pt, bodyoffsety=0pt, bodyverticalshift=0pt,
  roundedcorners=12, linewidth=2pt, titleinnersep=10pt, bodyinnersep=15pt
}{
  \draw[rounded corners=\blockroundedcorners, color=blocktitlebgcolor,
        fill=blockbodybgcolor, line width=\blocklinewidth]
    (blockbody.south west) rectangle (blocktitle.north east);
  \ifBlockHasTitle
    \draw[rounded corners=\blockroundedcorners,
          fill=blocktitlebgcolor, color=blocktitlebgcolor,
          line width=\blocklinewidth]
      (blocktitle.south west) rectangle (blocktitle.north east);
  \fi
}
\useblockstyle{MechValBlock}

\definetitlestyle{MechValTitle}{
  width=\paperwidth, roundedcorners=0, linewidth=0pt, innersep=20pt,
  titletotopverticalspace=15mm, titletoblockverticalspace=25mm
}{
  \draw[fill=titlebgcolor, color=titlebgcolor]
    (\titleposleft, \titleposbottom)
    rectangle (\titleposright, \titlepostop);
}
\usetitlestyle{MechValTitle}

\title{%
  \parbox{0.92\linewidth}{\centering\bfseries
    Mechanistic Validity\\[6pt]
    A Validity Theory, Research Methodology, and Evidence Standard\\[3pt]
    for Mechanistic Interpretability}}
\author{Elliot Tower\quad\texttt{elliot@elliottower.ai}}
\institute{}

\begin{document}

\maketitle[width=\paperwidth]

% ════════════════════════════════════════════════════════
% ROW 1
% ════════════════════════════════════════════════════════
\begin{columns}
\column{0.33}

\block{The Problem}{
  Mechanistic claims require five kinds of evidence---construct,
  measurement, internal, external, and interpretive validity---but
  researchers typically establish one or two and treat the claim as
  validated.

  \vspace{6mm}
  A strong causal result can conceal an unresolved construct definition.
  A reliable measurement can be measuring the wrong thing. A result
  established in one setting cannot automatically support broader
  generalization.

  \vspace{6mm}
  \textbf{Evidence along one dimension cannot substitute for another.}
  The field's dominant failure mode: establishing \emph{that} a
  component matters while leaving open \emph{what} it is.
}

\column{0.34}

\block{The Framework}{
  \textbf{36 criteria across 5 validity types.} Each claim receives a
  structured validity profile---pass, partial, or fail on each
  criterion. The tier is bounded by the weakest dimension, not the
  strongest.

  \vspace{5mm}
  \begin{center}
  \begin{tabular}{lcp{0.52\linewidth}}
    \toprule
    \textbf{Type} & \textbf{Criteria} & \textbf{Example question} \\
    \midrule
    Construct    & C1--C6  & Is the circuit a natural kind? \\
    Measurement  & M1--M7  & Does the metric measure what it claims? \\
    Internal     & I1--I12 & Is the causal claim warranted? \\
    External     & E1--E6  & Does the result generalize? \\
    Interpretive & V1--V5  & Is the interpretation justified? \\
    \bottomrule
  \end{tabular}
  \end{center}

  \vspace{4mm}
  The criteria form a dependency structure: the target mechanism must be
  coherent before it can be measured; the measurement must be
  trustworthy before it can support causal inference.
}

\column{0.33}

\block{Evidence Cards}{
  Each conclusion is tied to the specific measurement and criterion that
  support it. An evidence card records five fields:

  \vspace{5mm}
  \begin{enumerate}[leftmargin=*, itemsep=4mm]
    \item The \textbf{claim} under evaluation
    \item The \textbf{criterion} it addresses
    \item The \textbf{measurement} used to test it
    \item The \textbf{verdict} (confirmed, partial, inconclusive, disconfirmed)
    \item The resulting \textbf{tier}
  \end{enumerate}

  \vspace{5mm}
  Strength along one dimension stops standing in for weakness along
  another. A claim with strong causal evidence but no construct
  definition is recorded as such---the causal result is not discarded,
  but neither does it fill the construct gap.
}

\end{columns}

% ════════════════════════════════════════════════════════
% ROW 2
% ════════════════════════════════════════════════════════
\begin{columns}
\column{0.33}

\block{Case Study: The IOI Circuit}{
  Applied all 36 criteria to the indirect object identification circuit
  (Wang et al., 2023).

  \vspace{5mm}
  \begin{center}
  \begin{tabular}{lc}
    \toprule
    \textbf{Verdict} & \textbf{Count} \\
    \midrule
    \textcolor{confirmed}{\textbf{Confirmed}}            & 3  \\
    \textcolor{partial}{\textbf{Partially confirmed}}     & 18 \\
    \textcolor{inconclusive}{\textbf{Inconclusive}}       & 7  \\
    \textcolor{untested}{\textbf{Untested}}               & 5  \\
    \textcolor{disconfirmed}{\textbf{Disconfirmed}}       & 3  \\
    \bottomrule
  \end{tabular}
  \end{center}

  \vspace{4mm}
  The IOI circuit achieves descriptive adequacy and causal sufficiency
  but skips construct validity entirely. The analysis establishes
  \emph{that} specific heads matter without establishing \emph{what}
  functional role they fill as a natural kind.
}

\column{0.34}

\block{The Tier System}{
  Four tiers form a dependency chain. Each requires the evidence of all
  tiers below it.

  \vspace{5mm}
  \begin{enumerate}[leftmargin=*, itemsep=4mm]
    \item \textbf{Descriptive:} an observed pattern with no causal test
    \item \textbf{Causally Suggestive:} a single-method intervention
      supports the claim
    \item \textbf{Triangulated:} convergent evidence from independent
      methods with genuinely different failure modes
    \item \textbf{Mechanistically Adequate:} full validity profile with
      no disconfirmed criteria
  \end{enumerate}

  \vspace{5mm}
  Most published claims reach \textbf{Causally Suggestive}. The gap is
  between single-method causal evidence and convergent evidence from
  independent methods.
}

\column{0.33}

\block{Genetics Analogy}{
  The framework adapts four techniques from genetics that fill gaps
  interpretability's existing toolkit cannot cover:

  \vspace{5mm}
  \begin{enumerate}[leftmargin=*, itemsep=4mm]
    \item \textbf{Epistasis mapping:} do circuit components interact
      non-additively? Sub-additive interaction indicates a shared
      pathway; super-additive indicates mutual compensation.
    \item \textbf{Rescue experiments:} does restoring a single component
      recover function?
    \item \textbf{Sensitivity analysis:} does the choice of ablation
      method change the conclusion?
    \item \textbf{Complementation tests:} are the construct's labeled
      subdivisions functionally distinct? Adapts Benzer's cis-trans
      procedure.
  \end{enumerate}
}

\end{columns}

% ════════════════════════════════════════════════════════
% ROW 3
% ════════════════════════════════════════════════════════
\begin{columns}
\column{0.33}

\block{What This Shows}{
  Across the case study and surveyed literature, published claims
  establish \emph{that} a component matters and stop short of
  establishing \emph{what} it is.

  \vspace{6mm}
  \textbf{Construct validity is the systematic gap.} The field's
  interventions are sophisticated enough to support causal claims but
  the construct definitions those claims rest on are rarely tested.

  \vspace{6mm}
  Evidence cards make the gap visible: a claim with strong internal
  validity but untested construct validity is recorded at the tier its
  weakest dimension supports, and the missing evidence is named.
}

\column{0.34}

\block{What This Does Not Show}{
  \textbf{The framework evaluates claims, not methods.} A method is not
  invalid because a claim produced with it scores low---the claim may
  need additional evidence from complementary methods.

  \vspace{5mm}
  \textbf{The tier system is ordinal, not cardinal.} It orders the
  strength of evidence but does not quantify the distance between tiers.

  \vspace{5mm}
  \textbf{Single-paper validation is possible but incomplete.} A single
  paper can score high on many criteria; complete validation requires
  independent replication by a separate research group.

  \vspace{5mm}
  \textbf{The genetics analogies are tested on one circuit family.}
  Broader applicability awaits additional case studies.
}

\column{0.33}

\block{Takeaway}{
  \textbf{A mechanistic claim is warranted only to the extent that every
  required link between its concept, measurement, causal evidence,
  scope, and interpretation is supported.}

  \vspace{6mm}
  \begin{itemize}[leftmargin=*, itemsep=4mm]
    \item 36 criteria across 5 validity types
    \item Tier bounded by the weakest dimension
    \item Construct validity is the field's systematic gap
    \item Genetics provides four techniques with no current MI analogue
  \end{itemize}

  \vspace{8mm}
  \begin{center}
  \texttt{elliot@elliottower.ai}

  \vspace{3mm}
  {\small Paper and code:}

  \texttt{github.com/mechanistic-validity/}
  \end{center}
}

\end{columns}

\end{document}
